problem:  
Let \(G\) be a connected, simply connected eight‑dimensional **Hermitian Lie group** with Hermitian structure \((J,g)\) satisfying the **conformally balanced** condition  
\[
d(\Omega^{3}) = \theta \wedge \Omega^{3},
\]
where \(\Omega(X,Y)=g(JX,Y)\) and \(\theta\) is the Lee form.  
Assume \(G\) admits a **holomorphic conformal foliation** \(\mathcal{F}\) by complex surfaces (not assumed left‑invariant) whose leaves are **minimal** with respect to \(g\).  

Let \(D\subset TG\) denote the complex 2‑dimensional distribution tangent to the leaves. The distribution need not be left‑invariant.  

Investigate whether the conditions  
1. conformal balance, and  
2. minimality of the complex foliation,  
together imply that \(D\) must coincide (up to local automorphism of the complex Lie structure) with a **J‑invariant left‑invariant distribution**, equivalently that the foliation is locally generated by a **complex Lie subalgebra** of \(\mathfrak{g}\).  

If not, characterize obstructions to this Lie‑algebraic rigidity in terms of the torsion of the canonical connection associated with the SU(4)‑structure or of the curvature form of the Lee connection. In particular, determine whether vanishing of certain torsion components (for example, \(W_3\) or \(W_4\) types in the Gray–Hervella decomposition) enforces algebraic invariance of the foliation.  

Why is it a "good" problem:  
This problem asks whether analytic and variational geometric conditions—**conformal balance** and **minimality of holomorphic foliations**—are sufficiently strong to impose **algebraic rigidity** on the underlying Lie group. It removes the trivial case where the foliation is left‑invariant, turning the question into one about how curvature, torsion, and the Lee form of a conformally balanced Hermitian structure constrain the differential topology of complex foliations.  
It lies at the intersection of three active research directions: Hermitian curvature flows and SU(4)‑torsion theory, minimal and harmonic foliations in complex geometry, and algebraic structure theory of complex Lie groups. The statement is simple and natural yet open and potentially deep; resolving it would reveal how conformal‑analytic constraints shape the algebraic symmetries of Hermitian Lie groups, achieving the “narrow entrance, profound depth” characteristic of a good research problem.